%! Permutations and Transposes — Unit 2, Lesson 2.4 — Lesson Plan
\documentclass[10pt]{article}
\usepackage{linalg-article}
\usepackage{linalg-boxes}
\usepackage{graphicx}   % -article does NOT load graphicx; needed for thumbnails/screenshots
\graphicspath{{images/}}
\newcommand{\TallMath}[1]{$\displaystyle #1\rule[-1.4em]{0pt}{3.2em}$}
\newcommand{\xx}{\mathbf{x}}
\newcommand{\bb}{\mathbf{b}}
\newcommand{\cc}{\mathbf{c}}
\newcommand{\PP}{P}
\newcommand{\tp}{\mathsf{T}}
\newcommand{\UnitNumberName}{Unit 2: Solving Linear Equations Ax = b \quad}
\newcommand{\LessonNumberName}{Lesson 2.4: Permutations and Transposes}

\begin{document}
\begin{center}
    {\Large\bfseries \CourseName: \SchoolYear \\
    {\small\bfseries \UnitNumberName \LessonNumberName}}
\end{center}

\begin{tcolorbox}[colback=blush, colframe=burgundy, boxrule=0.9pt, arc=2mm,
    left=2mm, right=2mm, top=1.5mm, bottom=1.5mm]
    \textbf{Primary Objective:} Students can handle the one case Lesson~2.3 set aside --- a
    \emph{zero in the pivot}. They can build a \textbf{permutation matrix} $\PP$ (the identity with rows
    swapped) that records a row exchange, use it to fix a zero pivot, and state the resulting
    factorization $\PP A=LU$. They can also \textbf{transpose} a matrix $A\to A^{\tp}$ (turn its rows into
    columns), recognize a \textbf{symmetric} matrix ($A^{\tp}=A$), and apply the reversal rule
    $(AB)^{\tp}=B^{\tp}A^{\tp}$. They can explain \emph{why} a row swap never changes the solution to
    $A\xx=\bb$ --- it only reorders the equations.
\end{tcolorbox}

\renewcommand{\arraystretch}{1.4}
\begin{skillbox}[Priority Ideas \& Skills]{goldbox}
\begin{tabularx}{\linewidth}{@{}X|X@{}}
\textbf{Priority Ideas \& Skills} & \textbf{Key Understandings (the ``why'')} \\[2pt]
\hline
\vspace{-6pt}
\begin{itemize}[leftmargin=1.1em, nosep, after=\vspace{2pt}]
  \item Build a permutation matrix $\PP$ that swaps chosen rows; compute $\PP A$.
  \item Fix a zero pivot with a row swap: $\PP A=LU$.
  \item Transpose a matrix; recognize symmetric matrices ($A^{\tp}=A$).
  \item Use $(AB)^{\tp}=B^{\tp}A^{\tp}$ and $\PP^{\tp}=\PP^{-1}$.
\end{itemize}
&
\vspace{-6pt}
\begin{itemize}[leftmargin=1.1em, nosep, after=\vspace{2pt}]
  \item $\PP A$ just reorders the rows (equations) of $A$ --- so the solution $\xx$ is unchanged.
  \item 2.3 assumed every pivot was nonzero; the swap is how we handle the one case it isn't.
  \item Transpose reflects a matrix across its diagonal --- rows and columns trade places.
  \item The reversal rule is 2.2's socks-and-shoes again: undo/flip in the opposite order.
\end{itemize}
\end{tabularx}
\end{skillbox}

\begin{skillbox}[{Vocabulary, Concepts, \& Theorems}]{greenbox}
\renewcommand{\arraystretch}{1.3}
\begin{tabularx}{\linewidth}{@{}l X@{}}
\textbf{Permutation matrix $\PP$} & The identity matrix with its rows reordered. Multiplying $\PP A$ carries out that same reordering on the rows of $A$. \\
\textbf{Row exchange} & Swapping two rows (two equations). $\PP A=\bb$-side gets swapped too; the solution set is unchanged. \\
\textbf{$\PP A=LU$} & When a zero pivot blocks $A=LU$, swap rows first: $\PP A$ then factors into lower~$\times$~upper as usual. \\
\textbf{Transpose $A^{\tp}$} & Flip $A$ across its diagonal --- row $i$ becomes column $i$. Entry $(i,j)$ of $A$ moves to $(j,i)$. \\
\textbf{Symmetric matrix} & A square matrix with $A^{\tp}=A$ (mirror image across the diagonal). \\
\textbf{Reversal rule} & $(AB)^{\tp}=B^{\tp}A^{\tp}$ --- transpose flips a product's order (like the inverse). \\
\end{tabularx}
\end{skillbox}

\begin{fixedskillbox}[Activate Prior Knowledge \& Spiral Review]{blush}
    The Warm-Up rehearses the three ideas this lesson needs: \textbf{(1)} multiplying a matrix by a
    vector (Lesson~1.3) --- a swap matrix $\begin{bsmallmatrix} 0&1\\1&0 \end{bsmallmatrix}$ exchanges the
    entries, previewing what $\PP$ does to rows; \textbf{(2)} naming the first pivot (Lesson~2.1/2.3) ---
    students see a matrix whose top-left entry is $0$, so elimination \emph{cannot start}; and
    \textbf{(3)} reading an entry by its $(\text{row},\text{col})$ address (Lesson~1.3) --- the
    bookkeeping behind the transpose. Debrief item~2 aloud: this ``stuck'' pivot is exactly the gap 2.3
    left open, and today's row swap is the fix.
\end{fixedskillbox}

\begin{skillbox}[Hook]{blush}
    Put last lesson's method on a matrix that breaks it. Ask students to start elimination on
    $A=\begin{bsmallmatrix} 0&2\\1&3 \end{bsmallmatrix}$: the pivot would be the top-left entry, but it is
    $\mathbf{0}$ --- \textbf{``You can't divide by a zero pivot. Is this system broken, or can we
    rescue it?''} Draw out the fix everyone already knows from solving by hand: \emph{swap the two
    equations}. Swapping rows is itself a matrix --- $\PP=\begin{bsmallmatrix} 0&1\\1&0 \end{bsmallmatrix}$ ---
    and $\PP A=\begin{bsmallmatrix} 1&3\\0&2 \end{bsmallmatrix}$ now has a good pivot. That single idea,
    $\PP A=LU$, closes the last gap in solving $A\xx=\bb$.
\end{skillbox}

\begin{skillbox}[Lesson]{blush}
\begin{multicols}{2}
\textbf{1. A permutation matrix swaps rows.} Start from the identity and reorder its rows. The
$2\times2$ swap is $\PP=\begin{bsmallmatrix} 0&1\\1&0 \end{bsmallmatrix}$; then $\PP\begin{bsmallmatrix} a\\b \end{bsmallmatrix}=\begin{bsmallmatrix} b\\a \end{bsmallmatrix}$
and $\PP A$ swaps the two rows of $A$. Emphasize: $\PP$ only \emph{reorders} --- it never changes the
numbers, just their order.

\textbf{2. Fix a zero pivot: $\PP A=LU$.} On $A=\begin{bsmallmatrix} 0&2\\1&3 \end{bsmallmatrix}$ the pivot is
$0$, so swap: $\PP A=\begin{bsmallmatrix} 1&3\\0&2 \end{bsmallmatrix}$. Now the pivots are $1,2$ and it
factors normally (here it is already $U$, with $L=I$). For bigger matrices you swap first, then run the
ordinary elimination of 2.3 --- the multipliers fill $L$ as before.

\columnbreak

\textbf{3. Why the answer is safe.} Swapping two rows just reorders the equations. Swap the same two
entries of $\bb$ and the solution $\xx$ is \emph{identical} --- same equations, different order. So
solving $A\xx=\bb$ becomes solving $\PP A\xx=\PP\bb$, i.e.\ $LU\xx=\PP\bb$, by the two sweeps of 2.3.

\textbf{4. The transpose.} $A^{\tp}$ flips $A$ across its diagonal: row~$i$ becomes column~$i$, so the
$(i,j)$ entry lands at $(j,i)$. A \textbf{symmetric} matrix has $A^{\tp}=A$. Two rules matter:
$(AB)^{\tp}=B^{\tp}A^{\tp}$ (order reverses --- 2.2's socks-and-shoes) and $\PP^{\tp}=\PP^{-1}$ (a swap
is undone by the same swap).
\end{multicols}
\end{skillbox}

\begin{skillbox}[Active Monitoring]{redbox}
Circulate during the notes practice and Tier work. Watch for and correct:
\begin{itemize}[leftmargin=1.2em, nosep]
  \item swapping \emph{columns} instead of rows --- $\PP A$ (left multiply) swaps rows; only rows;
  \item forgetting to swap the matching entries of $\bb$ when solving $\PP A\xx=\PP\bb$;
  \item transposing by ``rotating'' --- it is a reflection across the diagonal, so the diagonal stays put;
  \item writing $(AB)^{\tp}=A^{\tp}B^{\tp}$ (wrong) instead of the reversed $B^{\tp}A^{\tp}$.
\end{itemize}
Cold-call prompts: ``What does $\PP$ do to the rows of $A$?'' ``Why doesn't the swap change $\xx$?''
``Where does the $(1,3)$ entry go under transpose?'' ``What makes a matrix symmetric?''
\end{skillbox}

\begin{skillbox}[Group Work \& Differentiation]{redbox}
\textbf{Swap, Factor, Flip} activity (tiered; mirrors the notes).
\begin{multicols}{3}
\textbf{Tier R --- Remediate.} Given a $2\times2$ $A$ with a zero pivot, write $\PP$, compute $\PP A$,
confirm the pivot is fixed, and solve $A\xx=\bb$.

\columnbreak
\textbf{Tier A --- Approaching.} Transpose a non-square matrix, test a matrix for symmetry, and
\emph{verify} $(AB)^{\tp}=B^{\tp}A^{\tp}$ on a $2\times2$ pair --- then interpret what symmetry means.

\columnbreak
\textbf{Tier E --- Extension.} A $3\times3$ with a zero pivot: write $\PP$, factor $\PP A=LU$, solve in
two sweeps, and justify why the row swap leaves the solution unchanged.
\end{multicols}
\end{skillbox}

\begin{skillbox}[Individual Work \& Assessment]{redbox}
\textbf{Exit Ticket} (independent, no notes): fix a $2\times2$ zero pivot with a permutation ($\PP$,
$\PP A$, solve), and transpose a matrix / test it for symmetry, plus a \textbf{conceptual/justification
check}: explain why swapping two equations (rows of $A$ and $\bb$) does not change the solution $\xx$.
Collect and sort into ``permutation \& transpose solid / swaps rows-vs-columns confused / shaky on why
$\xx$ is unchanged'' to decide whether the Unit~2 review opens with a quick re-teach.
\end{skillbox}

\begin{skillbox}[Reinforcement \& Extension]{goldbox}
\textbf{Homework:} build permutation matrices and compute $\PP A$ and $\PP^{\tp}\PP=I$; fix $2\times2$
zero pivots and solve; factor and solve a $3\times3$ $\PP A=LU$; transpose matrices and identify
symmetric ones; and verify $(AB)^{\tp}=B^{\tp}A^{\tp}$. \textbf{Extension:} show $A^{\tp}A$ is always
symmetric (using the transpose rules) and explain $\PP^{\tp}=\PP^{-1}$. \textbf{Preview:} Unit~2 wraps
up here --- next comes the unit test, then \emph{Unit~3, The Four Fundamental Subspaces}: now that we
can \emph{solve} $A\xx=\bb$, we ask which $\bb$ are reachable and what \emph{all} the solutions look
like.
\end{skillbox}
% ── Teacher notes (migrated from the component keys) ──
% One per component, titled for it. They live here rather than in the _key
% files so that every component is the same length blank and keyed.

\begin{teachernote}[Warm-Up]
Item~1 is a preview of what a permutation matrix does --- it reorders entries (and, applied on the left
to a matrix, reorders rows). Item~2 is the exact gap Lesson~2.3 left open: a zero pivot stalls
elimination, and today's row swap is the rescue. Item~3 sets up the transpose --- the $(1,2)$ and
$(2,1)$ entries ($1$ and $4$) are the two that trade places when we flip across the diagonal.
\end{teachernote}

\begin{teachernote}[Guided Notes]
The through-line: 2.3 assumed every pivot was nonzero; 2.4 supplies the missing case. \S1--\S2: a
permutation $\PP$ only \emph{reorders} rows, and a swap turns a stuck (zero) pivot into a good one, giving
$\PP A=LU$. Note that a $2\times2$ zero-pivot swap lands directly on $U$ with $L=I$ --- the nontrivial
$L$ shows up for $3\times3$ (activity/homework). \S3 is the conceptual payoff: a row swap reorders
equations, so with $\bb$ swapped too the solution is identical. \S4: transpose = reflect across the
diagonal (the diagonal stays put); symmetric means $A^{\tp}=A$; the reversal rule
$(AB)^{\tp}=B^{\tp}A^{\tp}$ mirrors 2.2's $(AB)^{-1}=B^{-1}A^{-1}$. Most common slips: swapping columns
instead of rows, and forgetting to reverse the order in the product rule.
\end{teachernote}

\begin{teachernote}[Group Activity]
Tier~R is the core skill (2$\times$2 swap $\to$ solve); watch for students swapping columns, or forgetting
to swap $\bb$. Tier~A is all transpose: reflect across the diagonal, test symmetry by comparing $(i,j)$
with $(j,i)$, and the reversal rule $(AB)^{\tp}=B^{\tp}A^{\tp}$ (order flips, exactly like the inverse in
2.2). Tier~E is the full $\PP A=LU$ payoff: swap first, then it factors and solves by 2.3's two sweeps ---
and the justification is the conceptual heart of the lesson. Note the $3\times3$ has a nontrivial $L$
(unlike the $2\times2$, where a swap lands straight on $U$).
\end{teachernote}

\begin{teachernote}[Exit Ticket]
Item~1 checks the whole permutation method: correct $\PP$, correct $\PP A$ and $\PP\bb$, and the solve
(watch for students who forget to swap $\bb$). Item~2 checks transpose and the symmetry test. Item~3 is
the conceptual target --- students must tie ``row swap'' to ``reordering equations'' and conclude the
solution is unchanged. Sort into ``permutation \& transpose solid / rows-vs-columns confused / shaky on
why $\xx$ is unchanged'' to plan the Unit~2 review.
\end{teachernote}

\begin{teachernote}[Homework]
Problem~3 is the full $\PP A=LU$ pipeline (swap $\to$ factor $\to$ two sweeps); note the nontrivial $L$
that a $3\times3$ produces. Problem~5 reinforces that transpose reverses a product's order (the same
socks-and-shoes structure as the inverse in 2.2), and the ``different matrix'' check makes the point
concrete. The extension previews why symmetric matrices $A^{\tp}A$ matter --- they return in Unit~4
(orthogonality) and Unit~6 (eigenvalues). Common slips: swapping columns instead of rows, forgetting to
swap $\bb$, and writing $A^{\tp}B^{\tp}$ instead of $B^{\tp}A^{\tp}$.
\end{teachernote}

\end{document}
